\subsection{Linearizing Common Nonlinearities --- Saturation, Friction, and Dead Zone}

Real physical systems rarely behave according to pure linear relationships. Every amplifier has a maximum output voltage, every mechanical joint exhibits friction that resists motion, and every gear train contains some amount of backlash that introduces play in the transmission. These nonlinear behaviors, while essential to understand for accurate system modeling, pose significant challenges when we attempt to apply the powerful methods of linear control theory. The technique of linearization, which we developed in the previous section, provides a systematic approach for approximating nonlinear system behavior near specific operating points. However, not all nonlinearities respond equally well to this treatment, and developing intuition for when linearization succeeds or fails is crucial for the practicing control engineer.

The fundamental question we must continually ask is whether the signals in our system will remain small enough to stay within the region where the linear approximation accurately represents the true nonlinear dynamics. When signals probe into strongly nonlinear regions, the linear model breaks down and our controller design may perform poorly or even become unstable. This section examines three of the most commonly encountered nonlinearities in control systems, deriving their linearized approximations from first principles and developing criteria for judging when these approximations remain valid.

The saturation nonlinearity appears whenever some physical quantity reaches a maximum or minimum value beyond which it cannot increase or decrease. Electrical amplifiers have power supply limits that constrain their output voltages, hydraulic valves cannot deliver flow beyond their rated capacity, and electric motors produce maximum torques at rated currents. Mathematically, we describe an ideal saturation nonlinearity through a piecewise relationship that passes the input through unchanged within some bounded range and clips the output at fixed limits when the input exceeds this range. For a saturation characteristic with lower limit $-y_{\max}$ and upper limit $+y_{\max}$, the input-output relationship takes the form shown in Table \ref{tab:nonlinearity_comparison}, which presents a side-by-side comparison of the three nonlinearities examined in this section.

\begin{table}[htbp]
\centering
\captionsetup{font=small}
caption{Comparison of the three common nonlinearities examined in this section, showing their mathematical definitions and linearization results.}
\label{tab:nonlinearity_comparison}
\begin{tabular}{|p{2.8cm}|p{3.8cm}|p{4.2cm}|}
\hline
\textbf{Nonlinearity} & \textbf{Nonlinear Relation} & \textbf{Linearized Model} \\
\hline
Saturation & $y = \begin{cases} y_{\max} & u > u_{\max} \\ u & |u| \leq u_{\max} \\ -y_{\max} & u < -u_{\max} \end{cases}$ & $y \approx K \cdot u$ with $K = \begin{cases} 0 & |u| \gg u_{\max} \\ 1 & |u| \ll u_{\max} \end{cases}$ \\
\hline
Coulomb Friction & $F = \begin{cases} -F_c \cdot \text{sgn}(\dot{x}) & \dot{x} \neq 0 \\ F_{\text{ext}} & \dot{x} = 0, |F_{\text{ext}}| < F_c \end{cases}$ & \makecell{Cannot linearize at $\dot{x}=0$;\\discontinuity at zero velocity} \\
\hline
Dead Zone & $y = \begin{cases} u - d & u > d \\ 0 & |u| \leq d \\ u + d & u < -d \end{cases}$ & $y \approx K \cdot u$ where $K = \begin{cases} 1 & |u| \gg d \\ 0 & |u| \ll d \end{cases}$ near origin \\
\hline
\end{tabular}
\end{table}

To derive the linearized model of saturation, consider an operating point defined by an input $u_0$ that produces an output $y_0$. The Jacobian of the saturation nonlinearity, which gives the local slope for linearization, equals $1$ everywhere within the linear region $|u| < u_{\max}$ and equals $0$ at all points where $|u| > u_{\max}$. At the exact boundary points $u = \pm u_{\max}$, the derivative does not exist because the function has a corner. For small perturbations $\Delta u$ around an operating point $u_0$ that lies well within the linear region, the linearized relationship becomes simply $y \approx u$, meaning the saturation nonlinearity contributes no gain modification and no phase shift to signals within its linear operating range.

However, the situation changes dramatically when the operating point approaches the saturation limits. If we attempt to linearize around an input $u_0$ that equals $90\%$ of the saturation limit, the linear model will predict that a small positive increase in input produces an equal increase in output. In reality, when the input reaches the saturation threshold, the output stops increasing entirely. The linear model predicts unbounded output growth while the actual system constrains output to the maximum value. This fundamental mismatch explains why control systems designed using linear models often fail spectacularly when actuators saturate, a phenomenon engineers call \textit{windup} that we shall examine in detail when we study PID control.

The practical implication of saturation linearization is that the linear model captures the system behavior accurately only for signals that remain sufficiently far from the saturation limits. A common rule of thumb suggests maintaining signals at less than $50\%$ of the saturation level to ensure the linear approximation remains adequate. When signals exceed this threshold, the designer must either redesign the system to provide additional actuator capability, implement anti-windup compensators that prevent the integrator states from growing without bound, or explicitly account for saturation in the controller design using techniques from the field of nonlinear control.

Coulomb friction represents a fundamentally different challenge for linearization because it introduces a discontinuity rather than a smooth nonlinearity. When two solid surfaces slide against each other, the friction force opposes the direction of motion with a magnitude that remains approximately constant regardless of sliding speed. This constant-magnitude resistance produces a force that switches sign instantaneously when the velocity changes sign, creating a piecewise-constant relationship that jumps discontinuously at zero velocity. The mathematical description of Coulomb friction states that the friction force $F_f$ equals $-F_c \cdot \text{sgn}(\dot{x})$ when the velocity $\dot{x}$ is nonzero, where $F_c$ is the Coulomb friction coefficient, and equals any value between $-F_c$ and $+F_c$ when the velocity is exactly zero (the system can remain stationary as long as the applied force does not exceed the static friction threshold).

The impossibility of linearizing Coulomb friction becomes apparent when we attempt to apply our Jacobian-based approach at zero velocity. Consider the operating point $\dot{x} = 0$. The right-hand derivative of friction force with respect to velocity, computed as the limit of the slope for small positive velocities, equals $-\infty$ because the friction force jumps from $0$ (at zero velocity) to $-F_c$ (for infinitesimal positive velocity). The left-hand derivative, computed for small negative velocities, equals $+\infty$ because the friction force jumps from $0$ to $+F_c$. Since these one-sided derivatives do not agree, no finite slope exists at zero velocity, and the linearization procedure fails completely.

This mathematical impossibility reflects a profound physical reality. The discontinuity in friction force means that the system behaves fundamentally differently at rest than in motion, and no linear differential equation can capture this switching behavior. When a control system commands a small torque to overcome friction and initiate motion, the actual response may differ dramatically from the linear prediction depending on whether the initial torque exceeds the static friction threshold. Engineers working with systems exhibiting significant Coulomb friction must either accept the limitation that linear models cannot capture stick-slip behavior and limit their designs to operating regimes where friction effects are negligible, or they must employ more sophisticated models that explicitly represent the friction discontinuity.

Dead zone nonlinearity occurs when a system fails to respond to inputs below some threshold magnitude. Mechanical transmissions contain backlash, the play between mated gear teeth that allows the driving gear to rotate a small amount before the driven gear begins to move. Hydraulic valves require a minimum pressure drop before flow begins. Electronic amplifiers exhibit threshold voltages below which no output appears. The dead zone characteristic shows zero output for inputs within the interval $[-d, d]$ and a linear relationship with slope unity for inputs beyond this interval. For positive inputs exceeding the dead zone threshold $d$, the output equals $u - d$, and for negative inputs below $-d$, the output equals $u + d$.

Linearizing the dead zone nonlinearity requires careful attention to the operating point. Consider an operating point $u_0$ that lies well outside the dead zone region, say $u_0 = 2d$. The local slope of the dead zone function at this point equals $1$, so the linearized relationship becomes $y \approx u - d$. This linear model predicts that the output responds one-to-one with the input, shifted by the constant offset $d$. For small perturbations around this operating point, the linear model accurately captures the system behavior because the perturbation range $|u - u_0| \ll d$ keeps the signal within the linear portion of the dead zone characteristic.

However, linearizing around the origin presents a different picture entirely. At $u_0 = 0$, the output $y_0 = 0$, but the slope of the dead zone function at the origin is zero for any small perturbation that remains within the dead zone. The Jacobian at this operating point equals $0$, producing the linearized model $y \approx 0$. This prediction states that the system produces no response to any small input around zero, which correctly captures the essential physics that inputs below the dead zone threshold produce no output. The linear model fails, however, when the input grows large enough to probe outside the dead zone region, because the actual system then begins responding while the linear model continues predicting zero output.

The following worked example demonstrates these concepts through a practical system that exhibits multiple nonlinearities simultaneously. Consider a position control system for a DC motor driving a lead screw to position a linear stage. The motor experiences Coulomb friction in its bearings, the power amplifier saturates at maximum current, and the mechanical coupling between the motor and lead screw exhibits backlash. We wish to design a proportional controller for this system and determine the operating conditions under which our linear design remains valid.

The motor dynamics without nonlinearities follow from Newton's second law applied to the rotational motion. The torque produced by the motor equals $J\ddot{\theta} + B\dot{\theta} + T_f(\dot{\theta}) + T_{\text{load}}$, where $J$ represents the moment of inertia, $B$ models viscous damping, $T_f$ captures Coulomb friction, and $T_{\text{load}}$ represents the load torque transmitted through the lead screw. The electrical dynamics of the motor and amplifier relate the applied voltage to the motor current through $L\frac{di}{dt} + Ri = V_{\text{applied}}$, where $L$ and $R$ represent the motor inductance and resistance respectively. For a permanent magnet DC motor, the torque constant $K_t$ relates current to torque through $T_{\text{motor}} = K_t i$, and the back-emf constant $K_e$ relates velocity to generated voltage through $V_{\text{back-emf}} = K_e \dot{\theta}$.

We linearize this system around an operating point defined by a constant reference position $\theta_0$ with zero velocity and zero acceleration. At this operating point, the motor current equals whatever value is necessary to balance the load torque plus friction, the back-emf equals zero, and the applied voltage equals the voltage drop across the resistance. The operating point current $i_0$ satisfies $K_t i_0 = T_{\text{load}} + T_f(0)$, where $T_f(0)$ lies somewhere between $-F_c$ and $+F_c$ depending on the history of motion.

Linearizing the electrical dynamics around this operating point yields the transfer function $V_{\text{applied}} - Ri_0 = (L s + R) I(s)$, where $I(s)$ represents the Laplace transform of the current deviation from $i_0$. Linearizing the mechanical dynamics gives $K_t I(s) = (J s^2 + B s) \Theta(s)$, where $\Theta(s)$ represents the angular position deviation from $\theta_0$. Combining these relations produces the open-loop transfer function from applied voltage deviation to angular position deviation:

\[
G(s) = \frac{\Theta(s)}{V_{\text{applied}}(s)} = \frac{K_t}{s(J s + B)(L s + R)}.
\]

This linear model accurately predicts system behavior as long as three conditions remain satisfied. First, the current deviation must remain small enough that the power amplifier does not saturate, which requires $|i_0 + i(t)| \leq I_{\max}$ for all time $t$. Second, the velocity must remain small enough that the motor does not enter the dead zone region of the amplifier input, which requires $|\dot{\theta}(t)| \geq v_{\min}$ for all time $t$ where $v_{\min}$ represents the minimum velocity for linear amplifier response. Third, the position error must remain small enough that the backlash region is never traversed, which requires $|\theta(t) - \theta_{\text{reference}}| < \theta_{\text{backlash}}$ for all time $t$ during motion through the dead zone region.

Each of these conditions places a constraint on the closed-loop system design. The saturation constraint limits the maximum torque the controller can request, which in turn limits the maximum acceleration the system can achieve and therefore constrains the achievable bandwidth of the position controller. The dead zone constraint on velocity effectively requires that the system never move so slowly that it enters the amplifier's nonlinear region, which proves problematic for fine positioning tasks where the system must move very slowly to achieve high positioning accuracy. The backlash constraint creates an effective dead zone in position control that produces limit cycling if the controller attempts to position precisely at a point within the backlash region.

The intuition we develop from studying these examples extends to the general question of when linearization succeeds or fails. Linearization succeeds when the signals in the system probe only the smooth portions of nonlinear functions and remain sufficiently far from discontinuities, saturation limits, or other strongly nonlinear regions. Linearization fails when signals grow large enough to encounter the corners, flat regions, or discontinuities that characterize strongly nonlinear behavior. The control engineer must always ask what signal amplitudes the closed-loop system will produce and whether these amplitudes remain within the linear operating range of all system components.

A useful mental model treats each nonlinearity as defining a region of validity in the state space of the system. The saturation nonlinearity defines a hypercube bounded by the maximum values each actuator can produce. The dead zone nonlinearity defines a central region where no response occurs, surrounded by linear response regions. Coulomb friction defines a switching surface at zero velocity that separates the stick regime from the slip regime. The linearized model accurately represents the true system only within the intersection of all these validity regions for the specific signals that occur in the closed-loop system.

In practice, engineers use several strategies to ensure their linear designs remain valid. Designing for conservative gain and phase margins provides robustness against the model errors introduced by linearization. Saturating controller outputs before they reach the actuators prevents the integrator windup problem. Implementing gain scheduling allows the controller parameters to change as the operating point moves through different regions of the nonlinear characteristic. When these strategies prove insufficient, the engineer must abandon linear methods entirely and employ nonlinear control techniques that explicitly account for the system nonlinearities.

The linearization techniques developed in this section and the previous one provide powerful tools for understanding and designing control systems. However, these tools come with inherent limitations that the engineer must always keep in mind. The next section extends our analysis to systems that cannot be linearized around a single operating point, introducing the concept of gain scheduling as a practical approach for controlling strongly nonlinear systems across their entire operating range.
